\exercisesheader{}

% 11

\eoce{\qtq{Apakah ini Bernoulli\label{is_it_bernouilli}} Tentukan apakah
setiap percobaan dapat dianggap sebagai percobaan Bernoulli yang saling
independen dalam situasi-situasi berikut.
\begin{parts}
\item Kartu-kartu yang dibagikan dalam satu tangan permainan poker.
\item Hasil setiap lemparan dadu.
\end{parts}
}{}

% 12

\eoce{\qt{Dengan dan tanpa pengembalian\label{with_without_replacement}}
Dalam situasi-situasi berikut, asumsikan bahwa separuh populasi yang
disebutkan adalah laki-laki dan separuh lainnya perempuan.
\begin{parts}
\item Misalkan Anda mengambil sampel dari sebuah ruangan yang berisi 10
orang. Berapa peluang terpilihnya dua perempuan berturut-turut jika
pengambilan dilakukan dengan pengembalian? Berapa peluangnya jika
pengambilan dilakukan tanpa pengembalian?
\item Sekarang, misalkan Anda mengambil sampel dari sebuah stadion yang
berisi 10.000 orang. Berapa peluang terpilihnya dua perempuan berturut-turut
jika pengambilan dilakukan dengan pengembalian? Berapa peluangnya jika
pengambilan dilakukan tanpa pengembalian?
\item Kita sering memperlakukan orang-orang yang diambil sebagai sampel
dari populasi besar sebagai saling independen. Berdasarkan temuan Anda pada
bagian~(a) dan~(b), jelaskan apakah asumsi ini masuk akal.
\end{parts}
}{}

% 13

\eoce{\qt{Warna mata, Bagian I\label{eye_color_geometric}} Sepasang suami
istri sama-sama bermata cokelat, tetapi membawa gen yang memungkinkan anak
mereka bermata cokelat (peluang 0.75), biru (0.125), atau hijau (0.125).
\begin{parts}
\item Berapa peluang bahwa anak pertama mereka yang bermata biru adalah
anak ketiga? Asumsikan bahwa warna mata setiap anak saling independen.
\item Secara rata-rata, berapa banyak anak yang akan dimiliki pasangan
tersebut hingga memperoleh anak bermata biru pertama? Berapa simpangan baku
banyaknya anak yang diperkirakan akan mereka miliki hingga anak bermata
biru pertama lahir?
\end{parts}
}{}

% 14

\eoce{\qt{Tingkat cacat\label{defective_rate}} Sebuah mesin yang memproduksi
transistor jenis khusus (salah satu komponen komputer) memiliki tingkat
cacat 2\%. Produksi dianggap sebagai proses acak dan setiap transistor
saling independen.
\begin{parts}
\item Berapa peluang bahwa transistor ke-$10$ yang diproduksi adalah
transistor cacat pertama?
\item Berapa peluang bahwa mesin tersebut tidak menghasilkan transistor
cacat dalam satu kelompok produksi yang terdiri atas 100 transistor?
\item Secara rata-rata, berapa banyak transistor yang diperkirakan akan
diproduksi hingga muncul transistor cacat pertama? Berapa simpangan bakunya?
\item Mesin lain yang juga memproduksi transistor memiliki tingkat cacat
5\%, dan setiap transistor diproduksi secara independen dari yang lain.
Secara rata-rata, berapa banyak transistor yang diperkirakan akan diproduksi
mesin ini hingga muncul transistor cacat pertama? Berapa simpangan bakunya?
\item Berdasarkan jawaban Anda pada bagian (c) dan (d), bagaimana peningkatan
peluang suatu peristiwa memengaruhi rata-rata dan simpangan baku waktu
tunggu hingga sukses?
\end{parts}
}{}

% 15

\eoce{\qt{Bernoulli, rata-rata\label{bernoulli_mean_derivation}}
Gunakan aturan peluang dari
Bagian~\ref{randomVariablesSection}
untuk menurunkan rata-rata variabel acak Bernoulli, yaitu variabel acak
$X$ yang bernilai 1 dengan peluang $p$ dan bernilai 0 dengan peluang
$1 - p$. Dengan kata lain, hitung nilai harapan variabel acak Bernoulli
umum.
}{}

% 16

\eoce{\qt{Bernoulli, simpangan baku\label{bernoulli_sd_derivation}}
Gunakan aturan peluang dari
Bagian~\ref{randomVariablesSection}
untuk menurunkan simpangan baku variabel acak Bernoulli, yaitu variabel
acak $X$ yang bernilai 1 dengan peluang $p$ dan bernilai 0 dengan peluang
$1 - p$. Dengan kata lain, hitung akar kuadrat varians variabel acak
Bernoulli umum.
}{}
